\section{Beilinson 谱投影子}\srcpage{35}

在 étale 理论中，Beilinson 谱投影子起关键作用。它给出一种从整个 $D(\Bun_G)$ 出发、以受控方式投影回性态良好的部分 $D_{\Nilp}(\Bun_G)$ 的办法。这有几个用途：\srcpage{36}

\begin{enumerate}
\item 它给出构造 $D_{\Nilp}(\Bun_G)$ 的对象的通用手段，特别地表明该范畴非平凡。值得注意的是，人们会把 Beilinson 谱投影子作用于 Whittaker 层，以构造 Langlands 函子。
\item 正如我们在下面所讨论的，它可用来产生 Hecke 特征层。
\item 计算 $D_{\Nilp}(\Bun_G)$ 上的 Frobenius 迹需要它。
\item 人们可以用它来证明定理 3.4 中的最终等价。目前 $D_{\Nilp}(\Bun_G)$ 指的是由定理 3.4 的性质 (1) 与 (2)（已知彼此等价）所定义的范畴。
\end{enumerate}

Arinkin、Gaitsgory、Kazhdan、Raskin、Rozenblyum 与 Varshavsky（2020b）中的原始构造，是一段关于范畴的复杂代数。我们将给出一个利用 Gestalt $\LocSys^{\et,\mathrm{Gest}}_{\Gd}$ 及其在整个 $D(\Bun_G)$ 上的作用的构造。这发生在基 Gestalt $[*]_D=\mathrm{Gest}(K_D)$ 之上：$K_D$ 是联系于六函子形式化 $D=\Ind D^{b}_{c}(-,\Ql)$（带 $(-_{*},-_{!})$-函子性）的对称幺半可表现 $(\infty,2)$-核范畴，$\mathrm{Gest}(K_D)$ 即其 Gestalt。

回忆注 5.7，我们有 $D(\Bun_G)\in K_D$，并带有 $\cO(\LocSys^{\et,\mathrm{Gest}}_{\Gd})_{1}=D_{\qc}(*/\Gd)_{\Ran}$ 的作用。此外，我们有子集
\[ i:\LocSys^{\et}_{\Gd}\times_{\mathrm{Gest}(D(\e))}[*]_D\subset\LocSys^{\et,\mathrm{Gest}}_{\Gd}. \]
我们于是可以把 $D(\Bun_G)$ 局部化到这个子集上——这是 $D(\Bun_G)$ 在 $D_{\qc}(\LocSys^{\et}_{\Gd})$ 作用于其上时的最佳逼近。因此，令
\[ D(\Bun_G)^{\mathrm{Hecke}}=D(\Bun_G)\otimes_{\cO(\LocSys^{\et,\mathrm{Gest}}_{\Gd})_{1}}\cO(\LocSys^{\et}_{\Gd})_{1}. \]

由于 $i$ 是 $1$-étale 且 $1$-proper（并且截断）的，它按 Scholze（2026，第 9 讲附录）容许 $!$-函子。因此，
\[ i^{*}:D(\Bun_G)\to D(\Bun_G)^{\mathrm{Hecke}} \]
容许一个截面 $i_{!}$，而 $D(\Bun_G)^{\mathrm{Hecke}}$ 是 $D(\Bun_G)$ 的收缩。另一方面，它在 $D(\Bun_G)$ 中的像必定包含于 Hecke-lisse 层之中，所以 $D(\Bun_G)^{\mathrm{Hecke}}$ 实际上是 $D_{\Nilp}(\Bun_G)\subset D(\Bun_G)$ 的收缩。\footnote{一个我们在此隐含使用的重要事实是：$D_{\Nilp}(\Bun_G)$ 是 restricted 的，即落在 $\PrL_{\e}\hookrightarrow K_D$ 的像中。} 由于后者已经是 $D_{\qc}(\LocSys^{\et}_{\Gd})$ 上的模，从而
\[ D(\Bun_G)^{\mathrm{Hecke}}=D_{\Nilp}(\Bun_G)\subset D(\Bun_G). \]

此外，函子 $i^{*}$ 现在定义 Beilinson 谱投影子
\[ \cP:D(\Bun_G)\to D_{\Nilp}(\Bun_G), \]
它是 $\cO(\LocSys^{\et,\mathrm{Gest}}_{\Gd})_{1}=D_{\qc}(*/\Gd)_{\Ran}$-线性的。

\begin{remark}
一个稍令人困惑之处是，Beilinson 谱投影子并不保证是包含
\[ i_{!}:D_{\Nilp}(\Bun_G)\subset D(\Bun_G) \]
的任何一侧的伴随。\srcpage{37}

人们应当想象 $i$ 既不是开浸入也不是闭浸入，而或许是局部闭浸入；在这种情形，$i_{!}$ 与 $i^{*}$ 不是伴随。
\end{remark}

作为紧生成范畴 $D(\Bun_G)$ 的收缩，形式上可知 $D_{\Nilp}(\Bun_G)$ 可对偶。实际上，可以证明它甚至是紧生成的：若使用与上面类似的手续，但把形式代数叠 $\LocSys^{\et}_{\Gd}$ 换成 Zariski 闭的代数子集，类似的映射 $i$ 将不再单射，但它将是 $0$-proper 的，因而函子 $i^{*}$ 有右伴随 $i_{*}$。这意味着在这一变体中 $i^{*}$ 保持紧对象，而这些紧对象给出 $D_{\Nilp}(\Bun_G)$ 的紧生成元。

关于对偶还可以说得更精确：在 $K_D$ 中，对象 $D(\Bun_G)$ 带有典范的自对偶数据，由余求值
\[ \Delta_{!}\Ql:\mathbf{1}\to D(\Bun_G)\otimes_{K_D}D(\Bun_G)=D(\Bun_G\times\Bun_G) \]
与求值
\[ \pi_{!}\Delta^{*}:D(\Bun_G)\otimes_{K_D}D(\Bun_G)\to\mathbf{1} \]
给出，即两者都由核 $\Delta_{!}\Ql$ 给出。这里 $\pi:\Bun_G\to*$ 表示投影，$\Delta$ 是对角。应用 Beilinson 谱投影子，可知 $D_{\Nilp}(\Bun_G)$ 也是自对偶的，相应的求值与余求值核由 Beilinson 谱投影子作用于 $\Bun_{G\times G}$ 上的 $\Delta_{!}\Ql$ 给出。\yizhu{原文此处作 $\Bun_{G\times G}$；按上下文，该核应是 $\Bun_G\times\Bun_G$ 上的 $\Delta_{!}\Ql$（前一式已如此写），此处照录原文。}{原文笔误（第 37 页；上下文推断）}

按 Gaitsgory 与 Raskin（2025a），包含
\[ D_{\Nilp}(\Bun_G)\subset D(\Bun_G) \]
保持紧对象；把它与 Arinkin、Gaitsgory、Kazhdan、Raskin、Rozenblyum 与 Varshavsky（2020b，第 4 节）结合起来——后者用 $(-^{!},-_{*})$-形式化中的伴随来刻画具有幂零奇异支撑的可构造层的性质——可以证明它实际上把紧对象送到 $(-_{*},-_{!})$-形式化中的 prim 层，从而给出一个等同
\[ D_{\Nilp}(\Bun_G)=\Ind D(\Bun_G)^{\prim}, \]
细节参见 Kuppel（2024）。我们注意到，在这一等同下，$D_{\Nilp}(\Bun_G)$ 的「神奇」自对偶不过是 prim 对偶——prim 对象经由取伴随，总带有反自等价。

借助 Beilinson 谱投影子，很容易定义 Hecke 特征层。即，对任意 $\Gd$-局部系统 $\varphi$，我们可以用 $D_{\qc}(\LocSys^{\et,\mathrm{Gest}}_{\Gd})$ 中由把 $i_{!}$ 作用于 $D_{\qc}(\LocSys^{\et}_{\Gd})$ 中的摩天大楼层而得到的对象来作用；这将投影到 $\varphi$ 对应的「Hecke 特征空间」上。遗憾的是，并不清楚这样能否得到非零对象。\srcpage{38}
